\section{Mass bridge: rest mass from self-confinement}
\label{sec:mass}

\subsection{Physical picture: wave trapped in a closed loop}
\label{subsec:mass-picture}

Rest mass in this framework is not acquired through the Higgs mechanism (spontaneous symmetry
breaking and a scalar condensate).
Instead, rest mass is identified with the total substrate energy of a wave mode that is permanently
trapped in a circularly closed, self-confined loop:
\begin{equation}
  mc^2 \;=\; E[{\vecX_{\mathrm{torus}}}]\Big|_{R=R^*},
  \label{eq:mass-energy}
\end{equation}
where $\vecX_{\mathrm{torus}}$ is the toroidal mode embedding and $R^*$ is the equilibrium major
radius.
The wave cannot disperse because it is topologically pinned (the $U(1)$ vortex carries
$\pi_1(U(1))=\mathbb{Z}$ winding) and geometrically self-confined (the geometric quartic prevents
radial expansion).
This is closer in spirit to the toroidal-electron proposal of Williamson and van der Mark
\citep{HuygensOptics2025}, which attributes the correct gyromagnetic ratio $g=2$ to a circularly
self-confined photon-like mode, than to the standard Higgs picture.
The substrate grounds that proposal in a concrete elastic medium.

\subsection{Complex structure from the time direction}
\label{subsec:complex-from-time}

The complex structure $\Psi\in\C^3$ used throughout this series is not an independent complex
field.
It arises from the carrier rotation of a real substrate field along the timelike worldtube axis.
For a band-isolated carrier at frequency $\omega_0$, the complex envelope is defined by
\begin{equation}
  \Psi \;\approx\; \xi + \frac{\ii}{\omega_0}\,\dot\xi,
  \label{eq:complex-from-time}
\end{equation}
where $\xi$ is the real lateral triplet and $\dot\xi$ is its time derivative.
The imaginary unit $\ii$ here is the quarter-period time shift (the time-evolution generator
$\ii\partial_t\Psi = H_{\mathrm{eff}}\Psi$); it is not an independent degree of freedom.
A single spacelike slice is a real snapshot and carries no phase; the $U(1)$ phase exists only
across the time extent of the worldtube.
Two adjacent time slices are the minimum data that fix the $U(1)$ --- exactly the structure of the
rotating-frame-periodic BVP solver that must be used for toroidal mode eigenstates.

This origin of the complex structure has a direct consequence for mass: the $U(1)$ carrier
frequency $\omega_0$ is the rate of rotation of the real field along the time axis, and the
energy stored in one carrier cycle is the substrate analogue of a quantum of energy.
The mass-frequency relation $E\approx\hbar\omega_0$ (in appropriate units) is therefore expected
to hold at the equilibrium amplitude, with $\hbar$ emerging from the substrate normalization
(\Cref{subsec:mass-normalization}).

\subsection{Toroidal mode energy and Derrick balance}
\label{subsec:toroidal-derrick}

A toroidal mode is a $U(1)$ vortex (line vortex in the trace sector) wrapped into a closed loop
of major radius $R$ and core radius $w$.
The total substrate energy decomposes into gradient and quartic contributions:
\begin{equation}
  E(R,A)
  \;=\; E_2(R,A) + E_4(R,A),
  \label{eq:torus-energy}
\end{equation}
where at leading order in $w/R$:
\begin{align}
  E_2(R,A) &\;\sim\; 2\pi R \cdot \varepsilon_2(A,w),
  \label{eq:torus-E2}\\
  E_4(R,A) &\;\sim\; 2\pi R \cdot \varepsilon_4(A,w),
  \label{eq:torus-E4}
\end{align}
with $\varepsilon_2$ and $\varepsilon_4$ the gradient and quartic energy densities per unit length
of the vortex core.
Both scale linearly with $R$ (circumference grows), so to leading order the total energy is linear
in $R$ --- this is the string-tension picture.
The correction that selects a finite equilibrium $R^*$ comes from the curvature energy of the
torus: turning the line vortex into a closed loop costs an extra curvature term $\propto 1/R$,
balancing the string-tension term $\propto R$:
\begin{equation}
  E(R) \;\approx\; T_{\mathrm{string}}\cdot 2\pi R + \frac{E_{\mathrm{curv}}}{R},
  \label{eq:torus-balance}
\end{equation}
with $T_{\mathrm{string}}$ the vortex string tension set by $\varepsilon_2+\varepsilon_4$ and
$E_{\mathrm{curv}}$ the curvature energy from bending the vortex line.
Setting $\mathrm{d}E/\mathrm{d}R=0$ gives the equilibrium radius
\begin{equation}
  R^* \;=\; \sqrt{\frac{E_{\mathrm{curv}}}{2\pi T_{\mathrm{string}}}}.
  \label{eq:torus-radius}
\end{equation}
The physical rest mass is then $mc^2 = E(R^*) = 2\cdot 2\pi R^* T_{\mathrm{string}}$:
the gradient and curvature contributions are equal at equilibrium.

\subsection{Time-link prestress as binding agent}
\label{subsec:time-link-binding}

The timelike link in the discrete 4D brane action (\Cref{eq:discrete-4d-action} of Paper~I)
carries a prestress $r_t = \alpha\beta\Delta t$ analogous to the spatial prestress parameter
$\alpha$.
Expanding the time-link norm term $-k_t r_t|\Delta R_t|$ produces a negative, lump-co-located
energy density in the carrier velocity field:
\begin{equation}
  \bar E_{t,\perp}(\rho)
  \;=\; -\frac{m\alpha}{4}\,(\omega A(\rho))^2 \;<\; 0,
  \label{eq:time-link-energy}
\end{equation}
where $\omega$ is the carrier frequency fixed by the worldtube closure condition
$\omega T = 2\pi n$ and $A(\rho)$ is the amplitude profile.
This induces a definite-negative longitudinal compression on the spatial links:
\begin{equation}
  \Delta u_\parallel^{\mathrm{ext}}(\rho)
  \;=\; -\frac{m\alpha}{4k_s b^2}\,(\omega A(\rho))^2
  \;\propto\; -\alpha\omega^2,
  \label{eq:time-link-compression}
\end{equation}
peaked at the vortex core $\rho\approx w$, zero on the axis.
This is externally sourced by the time-link prestress and carrier eigenvalue, not slaved to the
spatial relaxation.

Feeding \eqref{eq:time-link-compression} into the separable spatial cubic $V_3=+(k_s\alpha/2a^2)
\Delta u_\parallel(\xi_s^2+|\xi_\perp|^2)$ flips the net cross-coupling:
\begin{equation}
  g_{\mathrm{net}}
  \;=\; \frac{k_s\alpha}{4a^2}\,(1-\chi),
  \qquad
  \chi \;=\; 2\alpha\,\frac{\omega^2}{\omega_*^2}\,\frac{A_0^2}{a^2},
  \qquad
  \omega_* \equiv \frac{c_L}{b}.
  \label{eq:binding-condition}
\end{equation}
\emph{Binding} ($g_{\mathrm{net}}<0$, restoring force between $U(1)$ and $SU(3)$ sectors) holds
iff $\chi>1$.
This is reachable at the soliton sweet spot $\alpha\approx 0.5$--$0.8$, $\omega$ near the band
edge, $A_0/a\approx 0.3$.
The spatial-only verdict (net repulsive) is the $r_t\to 0$ limit; the time link genuinely flips
the sign.

\paragraph{Linchpin: closure-locked carrier frequency.}
The binding mechanism requires that $\omega$ be fixed externally by the worldtube closure
condition $\omega T=2\pi n$, not by spatial relaxation.
This is the case for the rotating-frame-periodic BVP solver (\texttt{PeriodicBC}/\texttt{solve\_block}):
the JFNK root-find relaxes node positions only; $\omega$ is not a degree of freedom.
Running the vortex through the breather eigensolver (where $\omega$ co-relaxes as a nonlinear
eigenvalue) would violate this condition and revert the binding verdict to net-repulsive.

\paragraph{Binding-gravity co-variation.}
The same $r_t$ (with $\alpha_t=\alpha$) sourcing the binding compression \eqref{eq:time-link-compression}
also sources gravitational time dilation (Paper~II): both coupling strengths scale as $\propto\alpha$.
As a result, binding strength and gravitational time-dilation strength co-vary in $\alpha$ and are not
independently tunable.
If a simulation can tune binding without tuning gravity, the identification $\alpha_t=\alpha$ is
falsified.

\subsection{Open normalization problem: Compton wavelength}
\label{subsec:mass-normalization}

The mass-energy identification \eqref{eq:mass-energy} requires three open normalization steps that
are not completed here:

\begin{enumerate}
  \item \textbf{Explicit toroidal energy integral.}
  Compute $E_2$, $E_4$, and $E_{\mathrm{curv}}$ in \eqref{eq:torus-balance} explicitly by
  integrating the substrate energy density over the torus geometry with the hedgehog-like core
  profile of \Cref{sec:matter}.

  \item \textbf{Equilibrium amplitude.}
  The amplitude $A$ is not fixed by the Derrick balance alone; it is selected by the energy
  normalization $E(R^*,A)=mc^2$ for a chosen target particle species.
  Express $A$ in terms of substrate parameters $(k_s,\alpha,a,m_{\mathrm{node}})$ and the
  target species mass.

  \item \textbf{Compton wavelength connection.}
  Show that $R^*=\hbar/mc$ using consistent derivations of $\hbar$ and $c$ from substrate
  parameters.
  Because $c^2=k_s a^2/m_{\mathrm{node}}$ and the natural action quantum is set by
  $k_s a^3$ (energy times volume) or $k_s a^2\cdot(a/c)$ (energy times time), the
  combination giving $\hbar$ in SI units must close the triangle
  $R^*\cdot mc = \hbar$ without additional free parameters.
\end{enumerate}

\subsection{Mass bridge open problems}
\label{subsec:mass-open}

\begin{enumerate}
  \item \textbf{Explicit toroidal energy and equilibrium radius.}
  Complete the calculation in \Cref{subsec:toroidal-derrick} for the substrate spring law
  (\Cref{eq:W4} of Paper~I).
  Identify $T_{\mathrm{string}}$ and $E_{\mathrm{curv}}$ in terms of $(k_s,\alpha,a,w,\omega_0)$.

  \item \textbf{Verify binding via closure-locked $\omega$ sweep.}
  Solve a stationary worldtube at each closure index $n$ (each imposed $\omega=2\pi n/(P\Delta t)$);
  on each converged state measure $\Delta u_\parallel(\rho)$ and the relaxed amplitude $A$.
  Confirm $\Delta u_\parallel\propto\omega^2$ across the converged family (dividing out $A$,
  which self-selects to the imposed $\omega$).

  \item \textbf{Winding-lock test.}
  Vary the topological winding $B$ at fixed amplitude on a converged baryon texture; check
  whether the P-odd holonomy $\gamma_\Gamma\propto B$ (soft winding-lock) or is flat in $B$
  (energetic-only).
  A $B$-slope of 1 confirms charge tracks winding; flat-in-$B$ indicates the energetic leg only.

  \item \textbf{Compton wavelength derivation.}
  Carry out the three-step normalization of \Cref{subsec:mass-normalization} and test whether
  $R^*=\hbar/mc$ closes without additional free parameters.
\end{enumerate}